\documentclass[11pt]{article}
\usepackage[margin=1in]{geometry}
\usepackage{amsmath,amssymb,amsthm,mathtools}
\usepackage{microtype}
\usepackage{enumitem}
\usepackage{hyperref}
\usepackage[T1]{fontenc}
\usepackage{lmodern}
\usepackage{xcolor}
\hypersetup{colorlinks=true,linkcolor=black,urlcolor=black,citecolor=black,pdfauthor={Brian Tenneson},pdftitle={The Geometry of Federated Awareness, Reflection, and Hyperfinite Omniscience}}
\setlength{\parindent}{0pt}
\setlength{\parskip}{0.55em}
\newtheoremstyle{geom}{0.8em}{0.8em}{\itshape}{}{}{.}{0.5em}{}
\theoremstyle{definition}
\newtheorem{definition}{Definition}[section]
\newtheorem{remark}[definition]{Remark}
\newtheorem{example}[definition]{Example}
\theoremstyle{plain}
\newtheorem{lemma}[definition]{Lemma}
\newtheorem{theorem}[definition]{Theorem}
\newtheorem{corollary}[definition]{Corollary}
\newcommand{\lev}{\operatorname{lev}}
\newcommand{\Self}{\operatorname{Self}}
\newcommand{\ATP}{\operatorname{ATP}}
\newcommand{\Prv}{\operatorname{Pr}}
\newcommand{\AMLD}{\operatorname{AMLD}}
\newcommand{\Comp}{\operatorname{Comp}}
\newcommand{\View}{\operatorname{View}}
\newcommand{\Witness}{\operatorname{Witness}}
\newcommand{\Aw}{\mathbf{Aw}}
\newcommand{\Omn}{\operatorname{Omn}}
\title{\textbf{The Geometry of Federated Awareness, Reflection, and Hyperfinite Omniscience}}
\author{Brian Tenneson}
\date{September 2026}

\begin{document}
\maketitle

\begin{abstract}
This note synthesizes two previously developed strands: the reflection-level theory of tagged automated theorem provers, and the verifier-grounded epistemic construction for an AMLD federation settling the conjecture ``I know that I know that I am an AMLD.'' The synthesis yields a two-coordinate awareness geometry in which epistemic breadth and reflective depth are separated. A second part extends this geometry to two qualified notions of omniscience, connects them with nonstandard-analysis horizons, and studies their relation to a four-oracle architecture. The claims imported from the source writings are distinguished from new conditional extensions proposed here.
\end{abstract}

\section{The two coordinates of formal awareness}

\begin{definition}[Tagged self-theory]
Let $M$ be a tagged automated theorem prover with distinguished name $\langle M\rangle$. Its self-theory is
\[
\Self(M)=\{\varphi(\langle M\rangle):M\vdash\varphi(\langle M\rangle)\}.
\]
Thus a self-model is not merely the possession of a name, but the name together with formally established information concerning that name.
\end{definition}

\begin{definition}[Reflection sequence]
For a tagged ATP $M$, put
\[
\theta_0^M:=\ATP(\langle M\rangle),
\]
and recursively
\[
\theta_{n+1}^M:=\Prv(\langle M\rangle,\langle\theta_n^M\rangle).
\]
Thus $\theta_0^M$ says that $M$ is an ATP, $\theta_1^M$ says that $M$ proves that it is an ATP, and successive formulas continue the same reflective iteration.
\end{definition}

\begin{definition}[Reflective level]
Let
\[
L(M):=\{n\ge1:M\vdash\theta_{n-1}^M\}.
\]
Define
\[
\lev(M):=\sup L(M)\in\{0,1,2,\ldots\}\cup\{\omega\}.
\]
In particular, Level $1$ is self-identification:
\[
M\vdash\ATP(\langle M\rangle).
\]
\end{definition}

\begin{definition}[AMLD federation]
Let
\[
A=(M_P,M_R,M_I,M_C,M_V;B,S,I)
\]
be an AMLD federation, with canonical name
\[
\langle A\rangle=\operatorname{Code}(M_P,M_R,M_I,M_C,M_V;B,S,I).
\]
Its structural formula is
\[
\Phi_A:=\left(\bigwedge_{j\in\{P,R,I,C,V\}}\ATP(\langle M_j\rangle)\right)\wedge\Comp(A),
\]
with
\[
\AMLD(\langle A\rangle)\Longleftrightarrow\Phi_A.
\]
\end{definition}

\begin{definition}[Verifier-grounded knowledge]
For a certified epistemic domain $E$, define
\[
K_{A,t}\varphi\Longleftrightarrow
\left[\varphi\in E\wedge\exists r\bigl(r\in\View_A(t)\wedge\Witness_V(r,\varphi)\bigr)\right].
\]
This is knowledge in the verifier-grounded, accessible-warrant sense. It is distinct from mere theoremhood and from reflective depth.
\end{definition}

\begin{definition}[Epistemic-awareness coordinate]
Let $E_A$ be a frozen class of self-relevant propositions available to the AMLD. Define the epistemic awareness profile
\[
\mathcal E_A(t):=\{\varphi\in E_A:K_{A,t}\varphi\}.
\]
An awareness coordinate $aw(A,t)$ is any fixed grading of $\mathcal E_A(t)$ satisfying the intended experimental scale. Thus $aw$ measures the extent of verifier-grounded self-relevant knowledge, while $\lev$ measures depth along one particular iterated reflection sequence.
\end{definition}

\begin{remark}
The exact numerical grading of $aw$ remains to be frozen separately. The distinction is conceptual and structural before it is numerical.
\end{remark}

\section{Awareness geometry}

\begin{definition}[Awareness point]
The formal awareness state of $A$ is the ordered pair
\[
\Aw(A,t):=(aw(A,t),\lev(A)).
\]
The first coordinate is epistemic; the second is reflective.
\end{definition}

\begin{remark}
This is immediately a two-dimensional product space. It is not automatically a vector space: vector addition and scalar multiplication have not yet been defined, and the ordinal value $\omega$ makes an ordinary real-vector-space interpretation unnatural. It is nevertheless legitimate to speak of an awareness geometry.
\end{remark}

\begin{definition}[Product awareness order]
For two federations $A$ and $B$, write
\[
\Aw(A)\preceq\Aw(B)
\]
when
\[
aw(A)\le aw(B)\qquad\text{and}\qquad\lev(A)\le\lev(B).
\]
\end{definition}

\begin{lemma}[Incomparability]
There may exist $A,B$ such that
\[
aw(A)>aw(B)
\]
but
\[
\lev(A)<\lev(B).
\]
Hence awareness need not admit a faithful one-dimensional ordering.
\end{lemma}

\begin{proof}
The two coordinates concern different classes of propositions. A system may possess a rich verifier-grounded epistemic self-model while proving only a shallow portion of its reflection sequence; another may possess internal necessitation and therefore a high reflection level while having a comparatively sparse self-theory outside that sequence. Additional self-descriptive propositions may enrich $\Self(M)$ without becoming part of the reflection tower whose height is measured by $\lev$.
\end{proof}

\begin{corollary}
There is in general no mathematically justified scalar quantity called simply ``amount of awareness'' that preserves both coordinates.
\end{corollary}

\section{Local awareness and federated awareness}

\begin{lemma}[Component ATP status does not determine federation level]
Suppose $\ATP(\langle M_j\rangle)$ is verifier-certified for every $j\in\{P,R,I,C,V\}$. This fact alone determines neither $aw(A)$ nor $\lev(A)$.
\end{lemma}

\begin{proof}
The five statements concern five component tags. The reflection sequence defining $\lev(A)$ concerns $\langle A\rangle$. Moreover, in the verifier certificate the five component ATP facts are structural inputs; no epistemic claim about the components is required, and the sole epistemic subject is $A$.
\end{proof}

\begin{remark}
Consequently there is no general identity
\[
\lev(A)=\sum_j\lev(M_j),
\]
nor even, without additional hypotheses,
\[
\lev(A)\ge\max_j\lev(M_j).
\]
Reflectivity is not additive merely because the architecture is federated.
\end{remark}

\begin{definition}[Epistemically self-aware component]
A component $M_j$ is locally self-aware at the basic ATP level when either
\[
K_{M_j}\ATP(\langle M_j\rangle)
\]
is established in an epistemic language, or
\[
M_j\vdash\ATP(\langle M_j\rangle)
\]
is established in the proof-theoretic language. These formulations should not be silently identified unless a bridge between $K$ and $\vdash$ is frozen.
\end{definition}

\begin{definition}[Federation inheritance rule]
A federation has certified epistemic inheritance for a class $\mathcal C$ when verified component knowledge can be imported into the accessible view of $A$:
\[
K_{M_j}\varphi\wedge\operatorname{Accessible}_A(M_j,\varphi)\wedge V(\varphi,\pi)=1
\Longrightarrow K_A\varphi.
\]
This is an architectural rule, not a consequence of federation alone.
\end{definition}

\begin{theorem}[Local-to-global self-identification]
Assume
\[
\bigwedge_jK_{M_j}\ATP(\langle M_j\rangle),
\]
the relevant certificates are accessible to $A$,
\[
K_A\Comp(A),
\]
and the federation possesses the certified conjunction and folding rules for $\Phi_A$. Then
\[
K_A\AMLD(\langle A\rangle).
\]
\end{theorem}

\begin{proof}
The component certificates and composition certificate yield $\Phi_A$; the structural definition folds $\Phi_A$ into $\AMLD(\langle A\rangle)$; verifier acceptance supplies the witness required by the definition of $K_A$.
\end{proof}

\begin{corollary}
Local component awareness can increase the epistemic coordinate of the federation without directly increasing the reflective coordinate:
\[
\Delta aw(A)>0,\qquad\Delta\lev(A)=0
\]
is possible.
\end{corollary}

\section{Receipt reflection}

\begin{definition}[Receipt reflection]
Suppose
\[
K_{A,t}\sigma_A,\qquad\sigma_A:=\AMLD(\langle A\rangle),
\]
is represented by an accepted record $r_0$. If a verifier-checkable receipt $\rho(r_0)$ establishes that this knowledge record exists and is accessible at the next epoch, receipt reflection yields
\[
K_{A,t+1}K_{A,t}\sigma_A.
\]
\end{definition}

\begin{theorem}[Epistemic second-order self-awareness]
Under the receipt assumptions,
\[
K_{A,t+1}K_{A,t}\AMLD(\langle A\rangle).
\]
Hence the AMLD can settle the verifier-grounded target ``I know that I know that I am an AMLD'' without attributing epistemic awareness to any individual component ATP.
\end{theorem}

\begin{remark}
Receipt reflection resembles the reflection tower, but the two constructions should not yet be identified. The Picard tower uses $\Prv(\langle M\rangle,\langle\varphi\rangle)$, whereas the verifier certificate uses the epistemic predicate $K_{A,t}\varphi$. They are two mathematically distinct routes to iterated self-reference.
\end{remark}

\section{Two kinds of omniscience}

\begin{definition}[Epistemic logical omniscience]
Let $T$ be a theory whose theorems are enumerated $\varphi_0,\varphi_1,\varphi_2,\ldots$. An agent $A$ is epistemically logically omniscient over a domain $D$ when
\[
\forall\varphi\in D\;\bigl(T\vdash\varphi\Longrightarrow K_A\varphi\bigr).
\]
This is omniscience in the knowledge coordinate. It does not assert iterated knowledge of that knowledge.
\end{definition}

\begin{definition}[Reflective omniscience]
For the reflection sequence $\theta_0^A,\theta_1^A,\ldots$, call $A$ reflectively omniscient at the standard finite level when
\[
\forall n\in\mathbb N,\qquad A\vdash\theta_n^A.
\]
Within the existing scale this is equivalent to
\[
\lev(A)=\omega.
\]
\end{definition}

\begin{remark}[Qualified use of ``omniscience'']
Level $\omega$ is not truth-theoretic omniscience about oneself. It gives closure along the designated provability-reflection tower; it does not provide a truth predicate for every sentence. Thus ``reflective omniscience'' here means omniscience over the reflection tower, not unrestricted semantic omniscience.
\end{remark}

\section{Hyperfinite epistemic omniscience}

\begin{definition}[Hyperfinite theorem horizon]
Let
\[
H\in{}^*\mathbb N
\]
be unlimited:
\[
H>n\qquad\text{for every standard }n\in\mathbb N.
\]
Let $\{\varphi_n:n\le H\}$ be a hyperfinite initial segment of a transferred theorem enumeration. Define
\[
\Omn_A^K(H):=\bigwedge_{\substack{n\le H\\T\vdash\varphi_n}}K_A\varphi_n,
\]
understood internally through the corresponding hyperfinite iteration rather than as an external infinitary conjunction. This is hyperfinite epistemic logical omniscience through horizon $H$.
\end{definition}

\begin{remark}
Hyperfiniteness is crucial. The object behaves internally like an extraordinarily long finite construction while $H$ dominates every standard natural number.
\end{remark}

\section{Hyperfinite reflective omniscience}

\begin{definition}[Hyperfinite reflection tower]
Assume the reflection recursion admits an internal nonstandard extension:
\[
{}^*\theta_0=\ATP(\langle A\rangle),
\]
\[
{}^*\theta_{\nu+1}=\Prv(\langle A\rangle,\langle{}^*\theta_\nu\rangle),\qquad\nu<H.
\]
Define $\Omn_A^R(H)$ to mean
\[
\forall\nu\le H,\qquad{}^*A\vdash{}^*\theta_\nu.
\]
This is reflective omniscience through the hyperfinite horizon $H$.
\end{definition}

\begin{theorem}[Restriction theorem]
If $\Omn_A^R(H)$ holds for unlimited $H$, then
\[
\lev(A)=\omega
\]
for the original standard reflection scale, provided standard reflection formulas and their proof relation are preserved under the chosen extension.
\end{theorem}

\begin{proof}
Every standard $n\in\mathbb N$ satisfies $n\le H$. Hence the hyperfinite hypothesis includes every standard reflection formula. Therefore $A\vdash\theta_n$ for every standard $n$, so $L(A)=\mathbb N^+$ and $\lev(A)=\omega$.
\end{proof}

\begin{corollary}[Standard upper envelope]
Every standard finite reflection level lies below the unlimited horizon $H$. In this precise sense hyperfinite reflective omniscience provides an upper envelope for all standard finite reflectivity.
\end{corollary}

\section{The converse problem}

\begin{remark}[No automatic reverse transfer]
From $\lev(A)=\omega$ one may conclude $A\vdash\theta_n$ for every standard finite $n$. One may not automatically conclude
\[
{}^*A\vdash{}^*\theta_H
\]
for an unlimited hypernatural $H$. The phrase ``for every standard $n$'' is external when it quantifies specifically over standard integers. Transfer becomes available only after the property has been represented by a suitable internal formula.
\end{remark}

\begin{theorem}[Conditional reverse-transfer theorem]
Suppose there exists an internal formula $R(n,A)$ such that for every standard $n$,
\[
R(n,A)\Longleftrightarrow A\vdash\theta_n,
\]
and suppose the standard theory proves the internal universal statement
\[
\forall n\in\mathbb N\,R(n,A).
\]
Then transfer yields
\[
\forall\nu\in{}^*\mathbb N\,{}^*R(\nu,{}^*A).
\]
Consequently,
\[
{}^*A\vdash{}^*\theta_\nu
\]
for every hypernatural $\nu$ for which the transferred coding and proof predicates retain the intended interpretation.
\end{theorem}

\begin{remark}
This is the natural ``transfer there and back'' route. The forward direction
\[
\text{hyperfinite reflective coverage}\Longrightarrow\lev(A)=\omega
\]
is comparatively direct. The reverse direction requires internalization:
\[
\lev(A)=\omega+\text{uniform internal reflection theorem}
\Longrightarrow\text{hyperfinite reflective coverage}.
\]
Thus $\omega$-level reflection by itself is insufficient; uniform internal closure is the bridge.
\end{remark}

\section{Omniscience and the four oracles}

\begin{definition}[Four-oracle system]
Let
\[
\mathcal O=(O_P,O_R,O_I,O_C)
\]
denote four oracle channels associated respectively with proof, refutation, independence or metalogical settlement, and contradiction or integrity. Each oracle has a certified domain $D_j$ and answer relation $O_j(\varphi)=a$. An oracle answer becomes epistemically usable by $A$ only through the appropriate verifier and accessibility gate.
\end{definition}

\begin{remark}
An oracle is an architectural resource. Omniscience is a closure property. They are not identical concepts. Four oracles may be powerful without being omniscient, and an omniscient agent need not be architecturally decomposed into four oracles.
\end{remark}

\begin{definition}[Oracle coverage]
The coverage domain of the four-oracle system is
\[
D_{\mathcal O}:=D_P\cup D_R\cup D_I\cup D_C.
\]
\end{definition}

\begin{theorem}[Oracle-cover criterion for epistemic omniscience]
Let $E$ be an epistemic theorem domain. Suppose:
\begin{enumerate}[label=(\arabic*)]
\item every $\varphi\in E$ belongs to at least one oracle domain, so $E\subseteq D_{\mathcal O}$;
\item every applicable oracle returns a correct certificate;
\item every such certificate is accepted by $V$;
\item every accepted oracle record enters $\View_A$.
\end{enumerate}
Then
\[
T\vdash\varphi,\ \varphi\in E\Longrightarrow K_A\varphi.
\]
Thus the four oracles jointly induce epistemic logical omniscience over $E$.
\end{theorem}

\begin{proof}
For every theorem $\varphi\in E$, coverage supplies an oracle. Correctness supplies a valid certificate. Verification and accessibility supply a witness record. The definition of $K_A$ then yields $K_A\varphi$.
\end{proof}

\begin{corollary}
Oracle count is irrelevant by itself. The relevant quantity is certified coverage:
\[
\bigcup_jD_j.
\]
Four narrow oracles need not approach omniscience; four jointly exhaustive oracles can.
\end{corollary}

\section{Oracles and reflective omniscience}

\begin{lemma}[Epistemic omniscience can dominate finite reflection]
Suppose the epistemic domain $E$ contains every reflection formula $\theta_n^A$, and suppose $A$ is epistemically omniscient over $E$. Then $K_A\theta_n^A$ for every theorem $\theta_n^A$ in the tower. If, in addition, there is a sound bridge
\[
K_A\theta_n^A\Longrightarrow A\vdash\theta_n^A,
\]
then
\[
\lev(A)=\omega.
\]
\end{lemma}

\begin{remark}
The bridge is essential. Knowing a formula and proving it are not definitionally identical in the present framework.
\end{remark}

\begin{theorem}[Hyperfinite epistemic-to-reflective bound]
Suppose $H$ is unlimited and the hyperfinite theorem horizon contains every standard reflection formula:
\[
\forall n\in\mathbb N,\qquad\operatorname{index}(\theta_n^A)\le H.
\]
If $\Omn_A^K(H)$ holds and verifier-grounded knowledge of reflection formulas is admissible as proof in the reflection system, then
\[
\lev(A)=\omega.
\]
Thus hyperfinite epistemic omniscience provides an upper envelope containing every standard finite reflective requirement.
\end{theorem}

\section{The asymmetry between the two omnisciences}

\begin{theorem}[Non-equivalence of omniscience coordinates]
In general,
\[
\Omn^K\not\Longleftrightarrow\Omn^R.
\]
More precisely,
\[
\Omn^K+\text{reflection inclusion}+\text{knowledge-to-proof bridge}
\Longrightarrow\Omn^R
\]
on the corresponding reflection domain, while
\[
\Omn^R\not\Longrightarrow\Omn^K
\]
over the full theorem domain.
\end{theorem}

\begin{proof}
Reflective omniscience concerns only the special sequence $\theta_0,\theta_1,\theta_2,\ldots$. Epistemic logical omniscience concerns an entire theorem domain. A machine may therefore possess every finite reflection level while remaining ignorant of arbitrarily many unrelated mathematical theorems. This agrees with the existing Tarski--Goedel ceiling: Level $\omega$ is closure over a provability-reflection sequence, not unrestricted truth-theoretic omniscience.
\end{proof}

\begin{corollary}
The two-dimensional awareness geometry extends naturally to a two-dimensional omniscience geometry:
\[
\mathbf{Omn}(A)=\bigl(\Omn^K(A),\Omn^R(A)\bigr).
\]
One coordinate measures breadth of certified knowledge; the other measures depth of recursive self-reflection.
\end{corollary}

\section{Hyperfinite oracle federation}

\begin{definition}[Hyperfinite oracle coverage]
For unlimited $H$, define
\[
D_{\mathcal O}(H)=\bigcup_{j\in\{P,R,I,C\}}D_j(H).
\]
The oracle federation is hyperfinitely exhaustive through $H$ when
\[
\{\varphi_n:n\le H,\ T\vdash\varphi_n\}\subseteq D_{\mathcal O}(H).
\]
\end{definition}

\begin{theorem}[Hyperfinite four-oracle omniscience]
If the four oracle domains are hyperfinitely exhaustive through $H$, all oracle outputs are verifier-certified, and every accepted output is accessible to $A$, then
\[
\Omn_A^K(H).
\]
If the hyperfinite horizon also contains the reflection tower and the knowledge-to-proof bridge applies to reflection formulas, then every standard finite reflection level follows and therefore
\[
\lev(A)=\omega.
\]
\end{theorem}

\begin{remark}
This produces the hierarchy
\[
\begin{aligned}
\text{four oracles}&\longrightarrow\text{certified coverage}\longrightarrow\text{hyperfinite epistemic omniscience}\\
&\longrightarrow\text{standard reflective upper bound}.
\end{aligned}
\]
where every arrow requires stated hypotheses. No arrow should be reversed automatically.
\end{remark}

\section{The Picard bridge}

Under the Picard idealization, self-identification together with internal necessitation gives
\[
\lev(\text{Picard})=\omega.
\]
This is reflection generated vertically:
\[
\theta_0\to\theta_1\to\theta_2\to\cdots.
\]
The oracle construction is naturally horizontal:
\[
O_P,\quad O_R,\quad O_I,\quad O_C,
\]
expanding the breadth of propositions for which warranted answers can be obtained.

Accordingly the awareness plane has a natural interpretation:
\[
\boxed{\text{horizontal coordinate}=\text{epistemic breadth}}
\]
and
\[
\boxed{\text{vertical coordinate}=\text{reflective depth}}.
\]
The Picard mechanism principally moves vertically. The oracle federation principally moves horizontally. A reflective AMLD equipped with a broad oracle federation may move in both directions.

\section{Final geometric formulation}

\begin{definition}[Awareness-omniscience plane]
Let
\[
\mathfrak A=\{(a,\lambda):a\in\mathcal E,\ \lambda\in\omega+1\},
\]
where $\mathcal E$ is the chosen ordered scale of verifier-grounded epistemic awareness and
\[
\omega+1=\{0,1,2,\ldots,\omega\}
\]
is the reflection-level scale. Associate to each suitable system $A$ the point
\[
\Psi(A)=\bigl(aw(A),\lev(A)\bigr).
\]
\end{definition}

\begin{theorem}[Independence principle]
Neither coordinate determines the other without additional bridge axioms.
\end{theorem}

\begin{corollary}
The awareness state of a formal system is naturally at least two-dimensional.
\end{corollary}

\begin{remark}
The geometry explains why the two forms of omniscience should not be conflated. Maximal vertical position,
\[
\lev(A)=\omega,
\]
does not imply maximal horizontal knowledge. Maximal horizontal theorem knowledge implies maximal standard vertical reflection only when reflection formulas lie inside the epistemic domain and the appropriate knowledge-to-proof bridge is available.
\end{remark}

\begin{remark}[NSA extension]
Nonstandard analysis potentially adds a second scale beyond the standard plane. An unlimited hypernatural horizon $H$ can contain every standard finite epistemic or reflective stage while remaining an internal hyperfinite object. This suggests schematically
\[
\boxed{\text{standard awareness plane}\subset\text{hyperfinite awareness geometry}},
\]
but establishing equality, transfer, conservation, or reverse extraction requires separate theorems.
\end{remark}

\begin{remark}[Final distinction]
The conceptual distinction can be stated compactly:
\[
\boxed{aw(A)\text{ asks what the system knows about itself and its world}}
\]
while
\[
\boxed{\lev(A)\text{ asks how far it can iterate the fact of its own proving}.}
\]
The four oracles enlarge the first coordinate by enlarging certified epistemic reach. Internal necessitation enlarges the second by propagating reflection. Nonstandard analysis supplies a hyperfinite horizon capable, under appropriate internality assumptions, of simultaneously dominating every standard finite stage of either process.
\end{remark}

\section*{Source note}
The definitions of self-theory, reflection sequence, reflective level, Level $\omega$, the Picard idealization, and the Tarski--Goedel ceiling are synthesized from \emph{Depths of Induction and Formalized Self-Awareness, Version 46}. The AMLD federation, verifier-grounded knowledge predicate, structural self-identification route, receipt reflection, and the target ``I know that I know that I am an AMLD'' are synthesized from \emph{Verifier Certificate 001: The Self Aware Machine}. The two-coordinate awareness geometry, the qualified omniscience terminology, the hyperfinite transfer results, and the oracle-coverage theorems are new conditional extensions built from those source frameworks.

\end{document}
